\begin{titlepage}
\thispagestyle{empty}
\begin{center}
{\Large Ontology of Continua}\\[0.6cm]
{\huge \textbf{Core 1.3}}\\[0.3cm]
{\Large Source-First Master Monograph}\\[1.2cm]

{\Large Alexander Yashin}\\[0.2cm]
{\normalsize Independent Researcher}\\[0.2cm]
{\normalsize ORCID 0009-0008-6166-0914}\\[1.0cm]

{\normalsize Version 1.3}\\[0.2cm]
{\normalsize April 13, 2026}\\[1.2cm]

\begin{minipage}{0.88\textwidth}
\small
\textbf{Publication role.} This volume is the canonical master monograph for OC Core 1.3.
It restores the full monograph-scale scientific shell, retains the inherited formal corpus of
Core 1.2 as the source-native substrate, and adds the explicit source-audit, theorem-fate,
chronology, and publication-boundary layers required for the 1.3 release.

\medskip

\textbf{Reader promise.} The goal of this document is not to compress the framework into a
short packet. The goal is to give the reader one inspectable scientific volume in which
definitions, theorem statements, proofs, figures, limitations, and provenance can be read in
their natural order without forcing the reader to reconstruct the theory from dashboard
artifacts or hidden companion ledgers.
\end{minipage}

\vfill

{\small Open foundation route: github.com/alexanderyashin/ontology-of-continua-core-main}\\
{\small Archival predecessor witness: Zenodo DOI 10.5281/zenodo.17903912}
\end{center}
\end{titlepage}

\begin{abstract}
Core 1.3 is the source-first revalidated master monograph of the Ontology of Continua.
It starts from the full formal shell and chapter tree that made Core 1.2 the first
axiomatically closed OC release, then subjects that inherited corpus to a stricter publication
discipline: explicit source witnessing, theorem-fate separation, chronology-aware repair notes,
nonclaim boundaries, and outward-ready scholarly packaging.

The result is intentionally larger than the short journal-facing article extracted from it.
This monograph is the didactic and formal master: it carries the structural theory, the
axiomatic and theorem-bearing backbone, the hierarchy of continua and embedding spaces, the
cross-level and domain modules, the falsifiability programs, and the explicit 1.3 materials
that explain what survives unchanged, what is revised, what is bounded, and what is deferred.

The scientific role of the volume is therefore dual. First, it preserves the monograph-level
readability that a serious foundational framework requires. Second, it supplies the canonical
source from which narrower journal-core and channel-specific artifacts can be derived without
silently widening or shrinking the defended claim. The monograph is thus both a scientific work
in its own right and the authoritative substrate for outward publication, review, translation,
and later projection into ESTRA, product, business, and public channels.
\end{abstract}

\clearpage
